\documentclass[german,10pt]{book}      \input{preambel}         \begin{document}  \setcounter{chapter}{2}


\chapter{Axiomatischer Zugang zur (Quanten-)Physik}
% Kap x
\label{chap_Axiom}

\info{Thomas Filk}{5.10.2025}%

In der Mathematik ist ein Axiomensystem ein minimaler Satz von Aussagen,\footnote{Ich setze hier
voraus, dass man entscheiden kann, ob eine Folge von Symbolen eine wohl definierte Aussage
bildet oder nicht. In der Mathematik wird dies im Rahmen der Symbollogik gekl\"art.} die
als wahr definiert werden und aus denen durch festgelegte Schlussfolgerungsregeln
weitere Aussagen als wahr oder falsch bewiesen werden k\"onnen. Die wesentlichen
Bedingungen an ein Axiomensystem sind\index{Axiom} 
\begin{itemize}
\item[(a)]
die Widerspruchsfreiheit: Es soll keine\index{Widerspruchsfreiheit}
Aussage geben, von der man beweisen kann, dass sie sowohl wahr als auch falsch ist, und 
\item[(b)] 
die Unabh\"angigkeit der Axiome, d.h., es soll nicht m\"oglich sein, f\"ur eines\index{Unabh\"angigkeit (von Axiomen)}
oder mehrere der Axiome unter Verwendung der anderen Axiome beweisen zu k\"onnen, 
dass sie richtig sind. 
\end{itemize}
Eine dritte Eigenschaft, die Vollst\"andigkeit, kann nach\index{Vollst\"andigkeit (von Axiomen)} 
den Unvollst\"andigkeitss\"atzen von G\"odel nicht immer garantiert werden.\index{Goedel@G\"odel}
\glqq Vollst\"andig\grqq\ bedeutet dabei, dass von jeder formulierbaren Aussage bewiesen
werden kann, ob sie richtig oder falsch ist. Dar\"uber hinaus spielt gelegentlich in der Praxis
eine Rolle, ob ein Axiomensystem reichhaltig bzw.\ interessant ist, d.h., ob sich viele
nicht selbstverst\"andliche Aussagen daraus ableiten lassen und ob diese Aussagen
zur Definition von interessanten Strukturen Anlass geben. 

Es gibt  viele Axiomensysteme in der Physik: die drei Newton'schen Postulate der Newton'schen
Mechanik, die Maxwell-Gleichungen zusammen mit der Lorentz-Kraft in der Elektrodynamik,
die Axiome der speziellen Relativit\"atstheorie (insbesondere die Konstanz der Lichtgeschwindigkeit
in allen Inertialsystemen unabh\"angig vom Bewegungszustand der Quelle) und der allgemeinen
Relativit\"atstheorie, die Haupts\"atze der Thermodynamik etc. Oft spricht man in der Physik
eher von Postulaten als von Axiomen, obwohl die Bedeutungen dieser beiden Terme kaum
zu trennen sind.\index{Postulate} 

Allgemein sind Postulate Aussagen, die man nicht mehr auf einfachere Aussagen
zur\"uckf\"uhren kann und die man als gegeben und richtig annimmt. In der Physik sollten
diese Aussagen und die aus ihnen ableitbaren Schlussfolgerungen nicht im Widerspruch
zum Experiment stehen. Oftmals handelt es sich sogar um Erfahrungstatsachen (z.B.\
die Konstanz der Lichtgeschwindigkeit in der SRT oder die Existenz einer Zustandsgr\"o\ss e
\glq Temperatur\grq\ in der Thermodynamik), die nicht unbedingt in jeder denkbaren Welt
gelten m\"ussten.   
Au\ss erdem sollen diese Aussagen untereinander nicht widerspr\"uchlich sein. 

In dieser Form k\"onnen Aussagen im Laufe der Zeit durchaus ihren Charakter als
Postulate verlieren: Aus heutiger Sicht kann man sagen, dass zu Keplers Zeit 
die Kepler'schen Gesetze den Charakter von Postulaten hatten -- aus der Empirie gewonnenen 
Regelm\"a\ss igkeiten, die sich nicht weiter begr\"unden lassen. Nachdem Newton seine
Gesetze formuliert hatte (die nun den Charakter von Postulaten hatten), konnte man die
Kepler'schen Gesetze aus den Newton'schen Gesetzen sowie dem Newton'schen 
Gravitationsgesetz ableiten. Das Newton'sche Gravitationsgesetz l\"asst sich wiederum aus den
Grundannahmen und Grundgleichungen der allgemeinen Relativit\"atstheorie ableiten.

Ganz allgemein wird man in der Physik zwei Arten von Postulaten unterscheiden: Zum einen
sind dies die Anforderungen, die man \"uberhaupt an eine physikalische Theorie bzw.\ ein
physikalisches Modell stellen muss, damit man diese Theorie als widerspruchsfreie Theorie
bezeichnen und in dieser Form akzeptieren kann. In Anlehnung an die Philosophie (und die
Kategorientheorie der Mathematik) k\"onnte man hier auch von Kategorien sprechen, die
erf\"ullt sein m\"ussen, damit wir \"uberhaupt von Physik sprechen k\"onnen. 
Zum anderen sind das Beobachtungstatsachen, die man nicht mehr weiter
erkl\"aren kann und daher einer Theorie als Postulate zugrunde legt. Zur ersten Gruppe kann
man beispielsweise das Postulat z\"ahlen, was \"uberhaupt (reine) physikalische
Zust\"ande sind. Zur zweiten Klasse
kann man das Postulat der Speziellen Relativit\"atstheorie z\"ahlen, dass sich Licht f\"ur
jeden inertialen Beobachter unabh\"angig vom Bewegungszustand der Quelle mit Lichtgeschwindigkeit
ausbreitet. 

In diesem Kapitel geht es haupts\"achlich um die erste Klasse von Postulaten, die allen
Theorien in der Physik gemein sind bzw.\ die jede Theorie spezifizieren muss, damit man sie
als Theorie bezeichnen kann. 

\section{Allgemeiner Formalismus}

In jeder physikalischen Theorie gibt es zwei Konzepte, auf denen die Theorie beruht: 
(1) Das Konzept der Observablen und (2) das Konzept der Zust\"ande.\index{Observable}\index{Zustand}
Ein axiomatischer Formalismus hat f\"ur eine Theorie zu kl\"aren, die diese beiden Konzepte mathematisch
repr\"asentiert werden und in welcher Beziehung sie zu den physikalischen Realisierungen stehen.

\subsection{Physikalische Realisierungen von Observablen und Zust\"anden}

In diesem Abschnitt werden die physikalischen -- man k\"onnte auch sagen \glqq experimentellen\grqq,
\glqq pragmatischen\grqq\ oder \glqq operationalen\grqq\ -- Realisierungen
einer Observablen und eines Zustands betrachtet. Diese sind unabh\"angig von einer Theorie und
gelten gleicherma\ss en in der Newton'schen Mechanik, der Quantenmechanik oder auch der
Thermodynamik. Die Unterschiede bestehen lediglich in der mathematischen
Repr\"asentation dieser Konzepte, die in den n\"achsten Abschnitten betrachtet werden.

Observable sind physikalische Gr\"o\ss en, die an einem physikalischen System gemessen werden
k\"onnen. Ihre Realisierung besteht in der Angabe eines Messprotokolls und einer Vorschrift,
wie man am Ende einer solchen Messung zu einer Zahl -- dem Ergebnis der Messung -- gelangt.
Es ist eine Erfahrungstatsache, welche Gr\"o\ss en das sind, und man kann nie sicher sein, dass
man alle an einem System beobachtbaren Gr\"o\ss en auch kennt.
John von Neumann definiert sogar ein physikalisches System durch die Menge der Observablen, 
die an diesem System beobachtet werden k\"onnen \cite{Neumann}. 

Die physikalische Realisierung eines Zustands erfolgt durch ein im Prinzip beliebig gro\ss es
Ensemble gleichartig pr\"aparierter Systeme. Gew\"ohnlich f\"uhren wir ein Experiment nicht nur
einmal durch sondern nach M\"oglichkeit unter denselben Ausgangsbedingungen gen\"ugend oft,
sodass wir statistische Auswerteverfahren verwenden k\"onnen. Die mathematische Vorstellung 
von einem Zustand als einem Erwartungswertfunktional auf der Menge der Observablen,
d.h.\ einer Abbildung ${\cal R} \longrightarrow {\rm Erw}({\cal R})$, wird dadurch realisiert, dass
man an einem ausreichen gro\ss en und repr\"asentativen Teilensemble der physikalischen Realisierung
eines Zustands die Messung einer Observablen ${\cal R}$ vornimmt und von den erzielten
Ergebnissen den Mittelwert bildet. 

Auf diese Weise k\"onnen wir auch
an demselben Zustand verschiedene Observablen messen, die sich an demselben System
nicht gleichzeitig bestimmen lassen (beispielsweise Ort und Impuls bei quantenmechanischen
Systemen oder allgemeiner zwei nicht kompatible Observable ${\cal R}$ und ${\cal S}$). 
Man unterteilt dazu das Ensemble gleichartig pr\"aparierter Systeme in
ebenfalls gro\ss e Teilensembles und misst an einem Teilensemble die Observable
${\cal R}$ und an einem anderen Teilensemble die Observable ${\cal S}$. Auf diese Weise erh\"alt man f\"ur
einen Zustand -- das Ensemble gleichartig pr\"aparierter Systeme -- die Erwartungswerte von
nicht kompatiblen Observablen in diesem Zustand.

Eine \"ahnliche Aussage gilt auch f\"ur Observable, die sich \textit{per definitionem} nicht durch
eine Einzelmessung bestimmen lassen, beispielsweise die Standardabweichung $\Delta {\cal R}$ einer
Observablen ${\cal R}$ oder ihren Mittelwert.


\subsection{Der axiomatische Rahmen einer Theorie}

Der axiomatische Rahmen einer Theorie hat zu kl\"aren, durch welche mathematische Strukturen
bzw.\ Objekte die physikalischen Observablen und Zust\"ande repr\"asentiert werden, und wie diese Strukturen
mit beobachtbaren Gr\"o\ss en -- das sind im Wesentlichen die Erwartungswerte von 
Observalben -- in Beziehung stehen. Das bedeutet, ein axiomatischer Rahmen hat die folgenden
Fragen zu kl\"aren:
\begin{enumerate}
\item
Durch welche mathematische Struktur werden Observable mathematisch dargestellt?
\item
Durch welche mathematische Struktur werden Zust\"ande mathematisch dargestellt?
\end{enumerate}
Neben der mathematischen Spezifizierung dieser beiden Konzepte muss ein allgemeiner 
Formalismus die folgenden Fragen beantworten k\"onnen: 
\begin{enumerate}
\setcounter{enumi}{2}
\item
Wie gelangt man aus der Kenntnis des Zustands zu Vorhersagen f\"ur die Resultate von 
Messungen einer Observablen an Systemen, die durch diesen Zustand beschrieben werden?
\item
Woher wei\ss\ man, durch welchen Zustand ein System zu beschreiben ist, 
wenn bestimmte Beobachtungsergebnisse vorliegen?
\item
Schlie\ss lich muss eine Theorie noch angeben, wie die
zeitliche Entwicklung des Systems beschrieben werden kann.
\end{enumerate}
Es gibt noch einige weitere Forderungen, die beispielsweise in der Standardformulierung
der Quantentheorie (hierunter verstehe ich eine Minimalform
der Kopenhagener Deutung, die im Wesentlichen die obigen Fragen beantwortet) 
teilweise schon nicht mehr erf\"ullt sind. Beispielsweise kann man von einer fundamentalen
Theorie fordern, dass sie f\"ur unsere Realit\"at eine Ontologie definiert, die auf das Universum
als Ganzes anwendbar ist. Dies ist f\"ur die Quantentheorie nicht der Fall.
Die folgenden Abschnitte kommentieren diese Konzepte.

\section{Observable und Zust\"ande}

Die beiden zentralen Begriffe -- Observable und Zust\"ande -- werden ausf\"uhrlich
in einem anderen Kurztext behandelt (siehe \href{../VDI_02_Observable_und_Zustaende}{Observable und Zust\"ande}). 
An dieser Stelle werden nur die
wichtigsten Eigenschaften im Zusammenhang mit diesen beiden Konzepten behandelt.

\subsection{Observable}

Ein allgemeiner physikalischer Formalismus muss f\"ur eine Theorie zun\"achst
kl\"aren, welche Observablen es gibt bzw.\ bekannt sind. Diese Observablen bilden jedoch nicht einfach 
nur eine strukturlose Menge, sondern es gibt auch 
Beziehungen zwischen den Observablen, die bekannt sein sollten und sich in einer mathematischen Darstellung
widerspiegeln m\"ussen.
In diesem Kapitel bezeichne ich die physikalische Observable
als Messvorschrift in Anlehnung an von Neumann \cite{Neumann} mit kaligraphischen Buchstaben, z.B.\
${\cal R}$, die zugeh\"orige mathematische Darstellung dieser Observablen kennzeichne ich 
durch $R$. 

Kennt man eine Observable ${\cal R}$, also die Vorschrift, wie man diese Observable an einem System 
messen kann, so kennt man auch die Observable $f({\cal R})$ f\"ur eine Funktion $f$: Man misst die
Observable ${\cal R}$ an einem physikalischen System und bildet von dem erhaltenen Messwert $r$ die
Funktion $f(r)$. Insbesondere ist die Observable $\alpha {\cal R}$ (f\"ur eine beliebige reelle Zahl $\alpha$)
durch die Messvorschrift von ${\cal R}$ gegeben, wobei jeder gemessene (an einer Zeigerstellung
abgelesene) Messwert mit $\alpha$ multipliziert wird. Lassen sich zwei Observable ${\cal R}$ und ${\cal S}$
gleichzeitig an einem System messen (das bedeutet, es gibt ein Protokoll, bei dem man gleichzeitig
einen Messwert sowohl f\"ur ${\cal R}$ als auch f\"ur ${\cal S}$ erh\"alt - diese Werte seien $r$ und $s$), 
dann ist 
die Observable $f({\cal R},{\cal S})$ definiert als die  Messvorschrift, bei der von den Messwerten die Funktion
$f(r,s)$ gebildet wird. 

Bei Observablen ${\cal R}$ und ${\cal S}$, die sich nicht gleichzeitig an einem System messen lassen -- solche
Observablen bezeichnet man als nicht kompatibel --, sind
allgemeine Funktionen $f({\cal R},{\cal S})$ zun\"achst nicht definiert. Insbesondere sind auch Ausdr\"ucke der Art
${\cal R}+{\cal S}$ oder ${\cal R}\cdot {\cal S}$ als Messvorschriften nicht definiert.   

Die mathematische Darstellung (Repr\"asentation) einer Observablen h\"angt von der Theorie
ab, die wir verwenden. In der klassischen Mechanik handelt es sich bei Observablen um 
Funktionen von Ort und Impuls (bzw.\ Geschwindigkeit), d.h.\ Funktionen auf dem
Phasenraum. In der Quantenmechanik
werden Observable durch sogenannte selbst-adjungierte bzw.\ hermitesche 
Operatoren\footnote{F\"ur unbeschr\"ankte Operatoren auf einem Hilbert-Raum besteht ein
Unterschied zwischen selbst-adjungiert und hermitesch, der hier jedoch nicht ber\"ucksichtigt
wird.} 
auf einem Hilbert-Raum -- einem Vektorraum mit einem Skalarprodukt -- dargestellt. 

Die mathematische Darstellung einer Observablen repr\"asentiert nicht den Messprozess, also
den dynamischen Vorgang der Messung, sondern eher die Informationen, die man durch solche
Messungen erlangen kann: die Menge der m\"oglichen Messwerte sowie die Zust\"ande, die mit
diesen Messwerten verbunden sind. 

\subsection{Zust\"ande}

Auch bei den Zust\"anden sollte man zwischen der physikalischen Realisierung und der
mathematischen Darstellung unterscheiden. Sehr oft definiert man einen Zustand als
ein \glqq Erwartungswertfunktional auf der Menge der Observablen\grqq. Diese zun\"achst
mathematische Definition kann man durchaus w\"ortlich interpretieren: Wenn wir ein System
durch einen Zustand beschreiben, erlaubt es dieser Zustand, jeder Observablen eine Zahl -- ihren
Erwartungswert in diesem Zustand -- zuzuordnen. Ein Zustand beschreibt also unsere Erwartungen
in Bezug auf die Ergebnisse, die bei einer Messung an einem System auftreten. 
Schr\"odinger bezeichnet einen Zustand als einen \glqq Katalog der Erwartung\grqq\ \cite{Schroedinger}. 

Dieser Katalog der Erwartung beruht auf unserem Wissen 
\"uber die Vergangenheit eines Sys\-tems bzw.\ eines Ensembles von Systemen. 
Im Wesentlichen bezieht sich dieses Wissen auf die Art, wie diese Systeme pr\"apariert wurden. 
Unsere Erwartungen in Bezug auf m\"ogliche Messergebnisse k\"onnen nur auf diesem Wissen
beruhen, und wenn uns bewusst ist, dass unser Wissen unvollst\"andig ist, verwenden
wir sogenannte gemischte Zust\"ande (in der klassischen Mechanik sind das Wahrscheinlichkeitsverteilungen
\"uber einem Zustandsraum, beispielsweise dem Phasenraum; in der Quantentheorie sind das
sogenannte Dichtematrizen). Unser Wissen \"uber die Vergangenheit
eines Sys\-tems muss nicht beliebig weit zur\"uckreichen: Es reicht das Wissen \"uber solche Aspekte, die
f\"ur die Vorhersagen zuk\"unftiger Messungen von Bedeutung sind. Das sind meist die letzten
Pr\"aparationen, die an einem System vorgenommen wurden.\footnote{Die Physik ist hier in einer
gl\"ucklichen Lage. M\"ochte man in der Psychologie den mentalen Zustand einer Person beschreiben,
k\"onnen beliebig weit zur\"uckliegende Ereignisse oder Erfahrungen f\"ur die Vorhersage des
zuk\"unftigen Verhaltens wesentlich sein.} Es gibt aber auch physikalische Systeme mit 
einem \glqq Ged\"achtnis\grqq, beispielsweise Spinglassysteme oder neuronale Netzwerke. 

Gelegentlich wollen wir auch einem Einzelsystem einen Zustand zuschreiben. An diesem
Punkt unterscheiden sich verschiedene Interpretationen, insbesondere die
sogenannten Ensemble-Interpretationen von der Kopenhagener Deutung. W\"ahrend
Ensemble-Interpretationen den Begriff der Wahrscheinlichkeit vermeiden und nur von
relativen H\"aufigkeiten sprechen -- die Born'sche Regel besagt dann, dass das
Absolutquadrat eines Skalarprodukts gleich der relativen H\"aufigkeit ist, mit der bei einer
Messung einer Observablen an einem Ensemble gleichartig pr\"aparierter Zust\"ande ein 
bestimmtes Resultat erzielt wird --
wendet die Kopenhagener Interpretation den Zustandsbegriff auch auf Einzelsysteme
an und spricht statt von relativer H\"aufigkeit von Wahrscheinlichkeit. Es gibt aber auch
Situationen, bei denen der Unterschied relevant wird: Beispielsweise gilt das sogenannte
No-Cloning Theorem, wonach ein unbekannter Zustand nicht dupliziert werden kann, 
nicht f\"ur Ensemblezust\"ande. An einem Ensemble kann der Zustand bestimmt und dann
beliebig oft kopiert werden. Diese Aussage gilt jedoch nicht f\"ur den Zustand, den wir einem
einzelnen System zuschreiben. 

Man bezeichnet einen Zustand -- realisiert als Ensemble -- als rein, wenn es keine Unterteilung 
dieses Ensembles in (f\"ur statistische Auswertungen ebenfalls gen\"ugend gro\ss e) 
Subensembles gibt, die f\"ur eines dieser Subensembles andere
Erwartungswerte liefert als das gesamte Ensemble. Diese Vorstellung von reinen Zust\"anden entspricht
auch unserer \"ublichen Vorstellung von reinen
Zust\"anden als maximales bzw.\ nicht mehr erweiterbares Wissen \"uber ein System. Entspricht
ein Ensemble von Systemen einem Gemisch, so l\"asst es sich in Subensembles aufteilen, die
weniger gemischten Zust\"anden (oder sogar reinen Zust\"anden) entsprechen. Reine Zust\"ande
lassen sich nicht mehr in noch reinere Zust\"ande aufteilen.
 
Mathematisch wird ein Zustand repr\"asentiert durch die Angabe einer Abbildung, die
jeder Observablen eine Zahl -- ihren Erwartungswert in diesem Zustand -- zuordnet. 
Ein Zustand ist somit eine mathematische Kodierung unseres Wissens \"uber die Art, wie ein
System pr\"apariert wurde, sodass wir dieses Wissen f\"ur die Vorhersage zuk\"unftiger
Messungen -- die Vorhersage eines Erwartungswerts f\"ur beliebige Observable -- an diesem 
System nutzen k\"onnen. 

Elementare Eigenschaften dieser Abbildung sind: 
\begin{enumerate}
\item 
Die Identit\"atsobservable, die immer nur
den Messwert 1 als Ergebnis liefert, soll den Erwartungswert 1 haben, 
\item 
eine positive Observable,
die immer nur positive Ergebnisse als Messwerte liefert, soll auch einen positiven Erwartungswert
haben, und 
\item
das $\lambda$-fache einer Observablen ${\cal R}$, die immer das $\lambda$-fache eines
Messwerts von einer Messung von ${\cal R}$ liefert, soll auch das $\lambda$-fache des Erwartungswerts
haben. 
\end{enumerate}
Weitere Eigenschaften h\"angen davon ab, welche Strukturen auf der Menge der
Observablen definiert sind. 

Oft unterscheidet man (insbesondere in der Quantentheorie) den Zustand von dem Erwartungswertfunktional, 
das dieser Zustand definiert. Der Zustand wird in der Quantentheorie z.B.\ durch einen normierten
Vektor $|\psi\rangle$ in einem Hilbert-Raum dargestellt, wohingegen die Abbildung 
${\cal R} \rightarrow {\rm Erw}({\cal R})$
durch ${\rm Erw}({\cal R})=\langle \psi|R|\psi\rangle$ gegeben ist. 

\subsection{Vorhersage und Pr\"aparation}

Der Begriff des Zustands als eine Abbildung, die jeder Observablen einen
Erwartungswert zuordnet, beinhaltet schon die Vorschrift, wie man aus einer Theorie
zu experimentell \"uberpr\"ufbaren Vorhersagen kommt. Wenn auf Seiten der Theorie
der Zustand eines Systems bekannt ist, kann man auch angeben, was man bei der
Messung einer bestimmten Observablen f\"ur einen Erwartungswert erh\"alt. Der
Erwartungswert ist dann der Mittelwert der Ergebnisse, die man bei der Messung 
einer Observablen an sehr vielen Systemen erh\"alt. 
Hierbei ist wichtig, dass der Zustand physikalisch durch ein Ensemble von Systemen
repr\"asentiert wird und nicht nur durch ein Einzelsystem. 

Diese Bedingung ist andererseits aber auch problematisch: Der Kosmos als Ganzes l\"asst
 sich nicht als Ensemble realisieren, das Gleiche gilt auch f\"ur komplexere Systeme wie
 beispielsweise einen menschlichen Organismus oder ein menschliches Gehirn. Zumindest
 reine Zust\"ande sind bei solchen Systemen nicht mehr realisierbar (das gilt schon f\"ur
 einfache thermodynamische Systeme). 
 
 Etwas problematischer ist die Frage, woher wir wissen, durch welchen Zustand ein 
 physikalisches System zu repr\"asentieren ist. Wir haben oben gesagt, dass wir die
 Vergangenheit eines Systems -- die Art wie es pr\"apariert wurde -- kennen m\"ussen. 
 Wir ben\"otigen somit eine Vorschrift, wie wir aus der Kenntnis der Pr\"aparation eines
 Systems den Zustand erhalten, durch den wir das System beschreiben. Diese Vorschrift
 ist oftmals ein eigenes Postulat (in der Quantentheorie beispielsweise das sogenannte
 Kollaps- oder Projektionspostulat).

\subsection{Die zeitliche Entwicklung eines Systems}

Die bisherigen Axiome sind allgemeiner Natur und enthalten noch keinerlei Aussagen
dar\"uber, wie sich ein physikalisches System im Verlauf der Zeit ver\"andert. 
Zun\"achst einmal muss \"uberhaupt gekl\"art werden, was \glqq zeitliche Entwicklung\grqq\
bedeutet. In fast allen physikalischen Theorien (eine Ausnahme bilden manche 
Modelle der Quantenkosmologie) wird die Zeit durch einen Parameter $t$ dargestellt
und als eindimensionales geordnetes Kontinuum gedacht. In solchen F\"allen bedeutet
\glqq zeitliche Entwicklung\grqq\ meist die Angabe einer Bewegungsgleichung, entweder f\"ur
die Zust\"ande (in der Quantenmechanik spricht man dann vom Schr\"odinger-Bild) oder
f\"ur die Observablen (das Heisenberg-Bild) eines Systems. In vielen F\"allen hat die
Bewegungsgleichung die Form einer Differentialgleichung. Es gibt aber auch Systeme,
bei denen es sich um eine Integralgleichung handelt -- meist Systeme mit einem Ged\"achtnis,
d.h.\ die zeitliche Entwicklung h\"angt m\"oglicherweise von Ereignissen ab, die in der
ferneren Vergangenheit des Systems stattgefunden haben.

\section{Die Postulate der Klassischen Mechanik}

In der klassischen Mechanik beschreiben wir einen (reinen) Zustand durch einen Punkt $(x,p)\in P$ im
Phasenraum $P$, also die Angabe eines Ortes $x$ und eines Impulses $p$ (oder, bei Vielteilchensystemen,
die Angabe aller Orte und aller Impulse). Die mathematische
Repr\"asentation einer Observablen ist eine
Funktion $F:P\rightarrow \mathbb{R}$ \"uber dem Phasenraum, also eine Funktion von Ort und Impuls.
Der Erwartungswert einer Observablen $F$ in einem Zustand $(x,p)$ ist einfach der Wert $F(x,p)$ dieser
Observablen an dem Punkt im Phasenraum, der den Zustand beschreibt. Dieser Wert ist eindeutig
und sollte (innerhalb der Fehlertoleranz der Messger\"ate) immer derselbe sein.  

Haben wir umgekehrt an einem System eine Observable ${\cal F}$ (mit zugeh\"origer mathematischer
Repr\"asentation $F$) gemessen und einen Wert $f$ erhalten, 
wissen wir, dass der Zustand ein Punkt im Phasenraum sein muss, der der Bedingung $F(x,p)=f$
gen\"ugen muss. Ein vollst\"andiger Satz von Observablen $\{{\cal F}_1,..., {\cal F}_n\}$ legt einen Punkt
im Phasenraum eindeutig fest. Misst man also diese Observablen und erh\"alt Messwerte
$\{ f_1,...,f_n\}$, so ist der Punkt $(x,p)$ durch die Bedingungen $\{F_1(x,p)=f_1,...,F_n(x,p)=f_n\}$
festgelegt und somit der reine Zustand des Systems bekannt. In einem minimalen vollst\"andigen
Satz von Observablen l\"asst sich auch keine Observable als Funktion der anderen Observablen
ausdr\"ucken. Hat der Phasenraum die Dimension $6N$ (f\"ur $N$ Punktteilchen in 3 Raumdimensionen),
so ben\"otigt man im Allgemeinen auch $6N$ Observable, um einen Zustand eindeutig festzulegen. 

Allgemein ist ein Zustand in der klassischen Mechanik eine Wahrscheinlichkeitsdichte $\omega(x,p)$ 
auf dem Phasenraum. Die Zuordnung eines Erwartungswerts f\"ur eine Observable mit mathematischer
Repr\"asentation $F$ ist dann
\begin{equation}
              \langle F \rangle_\omega = \int_{P} F(x,p) \omega(x,p)\,{\rm d}x\,{\rm d}p \, ,
\end{equation} 
wobei das Integral \"uber den gesamten Phasenraum $P$ zu nehmen ist. Die reinen Zust\"ande
erh\"alt man, indem man f\"ur $\omega$ speziell $\omega(x,p)=\delta(x-x_0) \delta(p-p_0)$
w\"ahlt. Der reine Zustand wird nun durch den Punkt $(x_0,p_0)$ im Phasenraum beschrieben.

Damit sind f\"ur die klassische Mechanik die ersten vier Punkte unseres allgemeinen Schemas
gekl\"art: (1) Observable sind Funktionen auf dem Phasenraum, (2) Zust\"ande sind Wahrscheinlichkeitsdichten
auf dem Phasenraum, (3) eine Wahrscheinlichkeitsdichte definiert zu einer Observablen einen Erwartungswert,
und schlie\ss lich (4) k\"onnen wir aus einem ausreichend gro\ss en Satz von Messwerten zu
(unabh\"angigen) Observablen -- beispielsweise Orts- und Impulsmessungen -- auch auf den Zustand 
schlie\ss en, durch den ein System zu beschreiben ist. Das letzte Axiom, die Zeitentwicklung, wird
dann beispielsweise durch die Newton'schen Axiome spezifiziert. Insbesondere das
zweite Newton'sche Axiom gibt an, wie man eine Bewegungsgleichung aufstellen kann, wenn die
Kr\"afte in einem System bekannt sind. Das erste (teilweise zusammen mit dem zweiten) Newton'schen Axiom 
gibt uns einen Hinweis, wie wir die Kr\"afte bestimmen k\"onnen, und das dritte Newton'sche Axiom
ist eine einschr\"ankende Bedingung, die von Kr\"aften erf\"ullt sein muss.   

\section{Die Postulate der Quantenmechanik}

Observable werden in der Quantenmechanik durch selbst-adjungierte bzw.\ hermitesche
Operatoren auf einem Hilbert-Raum dargestellt. Ein Zustand wird mathematisch durch einen Strahl
(einen eindimensionalen Unterraum) in diesem Hilbert-Raum beschrieben. Meist w\"ahlen wir zur einfacheren
Beschreibung einen auf eins normierten Vektor $|\psi\rangle$ auf diesem Strahl als Repr\"asentanten, wir
k\"onnen einen Zustand aber auch durch die Angabe eines eindimensionalen Projektionsoperators $P_\psi$
(der jeden Vektor in dem Hilbert-Raum auf diesen eindimensionalen Unterraum projiziert)
darstellen. In der Quantenmechanik von Punktteilen verwenden wir zur Beschreibung eines
Zustands oft eine sogenannte Wellenfunktion $\psi(x)$, die man aber als Vektor eines unendlich dimensionalen
Hilbert-Raums (dem Raum der quadratintegrierbaren Funktionen) auffassen kann. 
Das Erwartungswertfunktional, d.h.\ die Vorschrift, nach der
wir einer Observablen $A$ eine Zahl $\langle A \rangle_\psi$ -- ihren Erwartungswert in dem Zustand $\psi$ -- zuordnen, ist
dann 
\begin{equation}
       \langle A \rangle_\psi = \langle \psi | A |  \psi \rangle  = {\rm Spur}\,(P_\psi A) =
                \int_V \psi(x)^* A \psi(x)\,{\rm d} x \, .  
\end{equation}
Dies sind drei Darstellungen des Erwartungswertfunktionals, je nachdem, ob man einen Zustand durch
einen normierten Vektor, einen Projektionsoperator oder eine normierte Wellenfunktion \"uber einem Volumen 
$V$ repr\"asentiert. 

Das Kollapspostulat bzw.\ das von Neumann-L\"uders'sche Projektionspostulat gibt an, wie wir in der
Quantenmechanik ein System in einem bestimmten Zustand pr\"aparieren k\"onnen. Wurde eine Observable
${\cal R}$ an einem System gemessen und hat man den Messwert $r$ erhalten, so ist das System durch
einen Strahl zu beschreiben, der dem Eigenraum von $R$ zu dem Eigenwert $r$ entspricht. Ein 
vollst\"andiger Satz kompatibler Observabler legt durch ihre Messwerte diesen Eigenraum auf einen
eindimensionalen Strahl und damit einen reinen Zustand fest.

Ein allgemeiner Zustand wird in der Quantenmechanik durch eine Dichtematrix $\rho$
beschrieben und der Erwartungswert einer Observablen $A$ in einem Zustand $\rho$ ist
\begin{equation}
                  \langle A \rangle_\rho = {\rm Spur}\, \rho A \, .
\end{equation}
Eine Dichtematrix ist dabei durch folgende Bedingungen festgelegt: 
\begin{equation}
          \rho = \rho^\dagger ~(= (\rho^*)^T)  \hspace{1cm} 
          \rho \geq 0 ~( {\rm d.h.}~ \langle \psi |\rho| \psi \rangle \geq 0~~ \forall |\psi\rangle)  \hspace{1cm}
          {\rm Spur}\, \rho = 1 \, . 
\end{equation}
Jede Matrix mit diesen Eigenschaften l\"asst sich als Linearkombination von paarweise
orthogonalen Projektionsoperatoren $P_i$ schreiben:
\begin{equation}
             \rho = \sum_i p_i P_i \hspace{1cm}  {\rm mit}  ~~ \sum_i p_i = 1 ~~{\rm und} ~~p_i\geq 0 \, .
\end{equation}
Die orthogonalen Projektionsoperatoren $\{P_i\}$ beschreiben reine Zust\"ande 
und die Koeffizienten $\{p_i\}$ lassen sich als die Wahrscheinlichkeiten interpretieren, mit der
in dem gemischten Zustand $\rho$ der reine Zustand $P_i$ auftritt. 

Die Festlegung des Hamilton-Operators bestimmt dann die zeitliche Entwicklung eines
Sys\-tems, beispielsweise durch die Schr\"odinger-Gleichung oder die Heisenberg'schen
Bewegungsgleichungen.

\section{Die Postulate der Gleichgewichtsthermodynamik}

In der Gleichgewichtsthermodynamik sind die Observablen die makroskopischen physikalischen 
Gr\"o\ss en, die unter bestimmten Bedingungen kontrolliert bzw.\ konstant gehalten werden
k\"onnen. Neben der Energie $E$, der Teilchenzahl $N$ (eventuell den Teilchenzahlen $N_i$ von 
verschiedenen Substanzen oder Substanzen in verschiedenen Aggregatszust\"anden) und dem 
Volumen $V$ sind das beispielsweise auch die Temperatur $T$, die Entropie $S$, der Druck $p$ oder 
verschiedene chemische Potenziale $\mu_i$. 

Der sogenannte 0.te Hauptsatz der Thermodynamik
besagt, dass es eine beobachtbare Gr\"o\ss e $T$ gibt, genannt Temperatur, die die Eigenschaft
hat, dass sie f\"ur zwei Systeme, die eine ausreichend lange Zeit in Kontakt waren, gleich ist. Neben den
aus der Mechanik bekannten Gr\"o\ss en wie dem Volumen $V$ oder der Teilchenzahl $N$ wird
durch den 0.ten Hauptsatz f\"ur die Gleichgewichtsthermodynamik eine weitere Observable
postuliert, die sich nicht so einfach als Funktion \"uber dem Phasenraum der $N$ Teilchen
darstellen l\"asst. 

Zwischen diesen beobachtbaren (und kontrollierbaren) Gr\"o\ss en bestehen Beziehungen,
die f\"ur einen Gleichgewichtszustand charakteristisch sind. Beispielsweise l\"asst sich die Entropie
als Funktion der Energie, des Volumens und der Teilchenzahlen ausdr\"ucken. Das Gleiche
gilt f\"ur die Temperatur oder den Druck. Umgekehrt kann man aber auch beispielsweise
die Entropie als Funktion der Temperatur, des Volumens und der Energie schreiben. Je nach
unabh\"angigem Variablensatz, den man betrachtet, unterscheidet man verschiedene
Gesamtheiten. In der mikrokanonischen Gesamtheit betrachtet man meist die Energie als
thermodynamisches Potenzial und die Entropie, das Volumen und die Teilchenzahlen als
unabh\"angige Variable. Die Temperatur, der Druck und die chemischen Potenziale sind dann
abh\"angige Variable. In der kanonischen Gesamtheit ersetzt man die Entropie als unabh\"angige
Variable der mikrokanonischen Gesamtheit durch die Temperatur. Die Entropie wird dadurch
zu einer abh\"angigen Variablen. \"Ahnlich kann man auch das Volumen durch den Druck
als unabh\"angige Variable ersetzen, etc. 

Als Postulat fordert man meist, dass reine Gleichgewichtszust\"ande
eine st\"uckweise stetige Mannigfaltigkeit bilden. Ist der Gleichgewichtszustand als Punkt auf
dieser Mannigfaltigkeit bekannt, kann man alle Observablen als Funktion dieses Punktes
bestimmen. Dies ist \"ahnlich wie in der klassischen Mechanik. Der Unterschied zur klassischen
Mechanik besteht im Wesentlichen darin, dass die Menge der reinen Gleichgewichtszust\"ande
eine allgemeine Mannigfaltigkeit bildet. In der klassischen Mechanik ist diese Mannigfaltigkeit der
Phasenraum, d.h.\ ein (Ko)-Tangentialraum eines Konfigurationsraums, mit einer symplektischen
Struktur (der Poisson-Klammer f\"ur Observable). 




\begin{thebibliography}{99}
\bibitem{Neumann} von Neumann, John; \textit{Mathematische Grundlagen der
     Quantenmechanik}, Verlag von Julius Springer, Berlin 1935. 
\bibitem{Schroedinger} Schr\"odinger, E.; {\em Die gegenw\"artige
        Situation in der Quantenmechanik}; Die Naturwissenschaften 23
        (1935) 807--812, 823--828,  844-849.      
\end{thebibliography}

\end{document}

